\textbf{Время:} 01:15

\textbf{Тезис:} перестановка тождественных пионов даёт неразличимые варианты одного конечного исхода.

\textbf{На экране:} A и B слева → два конечных импульса справа → сплошные стрелки → пунктирные стрелки → вопрос справа.

\textbf{Текст:} Возьмём такую пару: справа указаны два положительных пиона. Слева точки A и B, а справа два зарегистрированных импульса.

Сначала пройдём по сплошным стрелкам: пион из A имеет импульс p-один, а пион из B — p-два. Теперь по пунктирным: импульсы поменялись местами.

В обоих вариантах зарегистрированы те же два импульса двух тождественных пионов. Перестановка не создаёт другого различимого конечного исхода. Поэтому вопрос справа — какой пион пришёл из A — не позволяет выбрать только один из этих вариантов. Квантовое описание должно учитывать оба.

\textbf{Переход:} Запишем для двух вариантов амплитуды и разберём, как из их суммы получается вероятность.

\textbf{Справка:} Для бозонов используется симметричная сумма амплитуд. Это принцип квантового описания тождественных частиц, а не просто недостаточная точность детектора. Для пары положительного и отрицательного пиона такого обменного члена нет; взаимодействия пары всё равно могут создавать корреляции.

\textbf{Источники:} Lisa et al. (2005), рис. 1 и §2.1.
